\documentclass[11pt,a4paper]{article}

% ---------------------------------------------------------------- langue & fontes
\usepackage{fontspec}
\usepackage[french]{babel}
\frenchsetup{SmallCapsFigTabCaptions=false,ThinSpaceInFrenchNumbers=true}
\setmainfont{LibertinusSerif-Regular.otf}[
  BoldFont=LibertinusSerif-Bold.otf, ItalicFont=LibertinusSerif-Italic.otf,
  BoldItalicFont=LibertinusSerif-BoldItalic.otf, Numbers=OldStyle]
\setsansfont{Inter-Regular.otf}[
  BoldFont=Inter-Bold.otf, ItalicFont=Inter-Italic.otf,
  BoldItalicFont=Inter-BoldItalic.otf, Scale=0.94]
\setmonofont{FiraMono-Regular.otf}[BoldFont=FiraMono-Bold.otf, Scale=0.84]
\newfontfamily\tabfont{Inter-Regular.otf}[BoldFont=Inter-Bold.otf,
  ItalicFont=Inter-Italic.otf, BoldItalicFont=Inter-BoldItalic.otf,
  Scale=0.94, Numbers=Lining]

% ---------------------------------------------------------------- mise en page
\usepackage[a4paper,top=2.5cm,bottom=2.3cm,left=2.1cm,right=2.1cm,
            headheight=17pt,headsep=13pt,footskip=20pt]{geometry}
\usepackage{microtype}
\usepackage{xcolor}
\usepackage{booktabs,tabularx,longtable,array,colortbl,multirow}
\usepackage{titlesec,titletoc}
\usepackage{fancyhdr}
\usepackage{enumitem}
\usepackage[most]{tcolorbox}
\usepackage{fancyvrb}
\usepackage{pifont}
\usepackage{fontawesome5}
\usepackage{needspace}
\usepackage{ragged2e}
\usepackage[hidelinks]{hyperref}

% ---------------------------------------------------------------- palette
\definecolor{encre}   {HTML}{16181D}
\definecolor{accent}  {HTML}{1F4E5F}
\definecolor{accent2} {HTML}{8C4A2F}
\definecolor{gris}    {HTML}{6B7280}
\definecolor{filet}   {HTML}{D5DADE}
\definecolor{fondclair}{HTML}{F3F6F7}
\definecolor{fondcode}{HTML}{F7F7F4}
\definecolor{vert}    {HTML}{17693F}
\definecolor{rouge}   {HTML}{9B2C2C}
\definecolor{ambre}   {HTML}{A86A12}
\color{encre}

\hypersetup{colorlinks=true,linkcolor=accent,urlcolor=accent,
  pdftitle={Reunion de cadrage - Projet NLP 2},
  pdfsubject={Jalon J2 - choix des cas d'usage, cadrage, outils novateurs},
  pdfauthor={Equipe projet NLP 2 - M2 Data et IA, FGES}}

% ---------------------------------------------------------------- symboles
\newcommand{\okmark}   {\textcolor{vert}{\ding{51}}}
\newcommand{\komark}   {\textcolor{rouge}{\ding{55}}}
\newcommand{\warnmark} {\textcolor{ambre}{\faExclamationTriangle}}
\newcommand{\checkbox} {\textcolor{gris}{\ding{113}}}
\newcommand{\cutmark}  {\textcolor{accent2}{\ding{34}}}
\newcommand{\starmark} {\textcolor{ambre}{\ding{72}}}
\newcommand{\wipmark}  {\textcolor{ambre}{\faTools}}
\newcommand{\code}[1]  {{\ttfamily #1}}
\newcommand{\rulebreak}{}

% ---------------------------------------------------------------- titres
\titleformat{\section}[hang]
  {\sffamily\bfseries\LARGE\color{accent}}{\thesection}{13pt}{}
  [\vspace{2pt}{\color{filet}\titlerule[0.9pt]}]
\titlespacing*{\section}{0pt}{24pt plus 4pt minus 2pt}{11pt}
\titleformat{\subsection}
  {\sffamily\bfseries\large\color{encre}}{\thesubsection}{9pt}{}
\titlespacing*{\subsection}{0pt}{16pt plus 3pt minus 1pt}{6pt}
\titleformat{\subsubsection}
  {\sffamily\bfseries\normalsize\color{accent2}}{}{0pt}{}
\titlespacing*{\subsubsection}{0pt}{12pt plus 2pt minus 1pt}{4pt}
\setcounter{secnumdepth}{2}
\setcounter{tocdepth}{2}

% ---------------------------------------------------------------- en-têtes
\pagestyle{fancy}\fancyhf{}
% en-tête courant, redéfini par build.py selon le document
\newcommand{\entetecourant}{Projet NLP 2}
\fancyhead[L]{\sffamily\scriptsize\color{gris}\entetecourant}
\fancyhead[R]{\sffamily\scriptsize\color{gris}\thepage}
\renewcommand{\headrulewidth}{0pt}
\renewcommand{\headrule}{\vspace{-4pt}{\color{filet}\rule{\headwidth}{0.4pt}}}
\fancypagestyle{plain}{\fancyhf{}\renewcommand{\headrule}{}%
  \fancyfoot[C]{\sffamily\scriptsize\color{gris}\thepage}}

% ---------------------------------------------------------------- paragraphes & listes
\setlength{\parindent}{0pt}
\setlength{\parskip}{5.5pt plus 1.5pt minus 0.5pt}
\linespread{1.05}
\raggedbottom
\widowpenalty=10000 \clubpenalty=10000
\setlist[itemize]{leftmargin=1.15em,itemsep=2pt,topsep=3pt,parsep=0pt,label=\textcolor{accent}{\textbullet}}
\setlist[enumerate]{leftmargin=1.5em,itemsep=2.5pt,topsep=3pt,parsep=0pt,
  label=\textcolor{accent}{\sffamily\bfseries\arabic*.}}

% ---------------------------------------------------------------- tableaux
% \hspace{0pt} : glue initiale, sans quoi TeX refuse de couper le premier
% mot de la cellule — cause de débordements en colonne étroite.
\newcolumntype{L}[1]{>{\raggedright\arraybackslash
  \rightskip=0pt plus 1fil\relax\hspace{0pt}}p{#1}}
% \hspace{0pt} : sans cette glue, le whatsit de \color empêche TeX
% de couper le premier mot de la cellule (débordements en en-tête).
\newcommand{\thead}[1]{{\tabfont\bfseries\color{accent}\hspace{0pt}#1}}
\newenvironment{tabletype}
  {\par\addvspace{7pt}\begingroup
   \tabfont\setlength{\tabcolsep}{4pt}\renewcommand{\arraystretch}{1.26}
   \setlength{\LTleft}{0pt}\setlength{\LTright}{0pt}
   \setlength{\extrarowheight}{1pt}
   % élasticité infinie à droite : TeX ne coupe plus un mot pour « remplir »,
   % mais coupe encore lorsque c'est le seul moyen d'éviter un débordement.
   % césure interdite : les largeurs minimales sont mesurées, chaque mot tient.
   % Un débordement, s'il en restait un, serait signalé par le journal.
   \hyphenpenalty=9000\exhyphenpenalty=9000
   \emergencystretch=0pt\hbadness=10000\relax}
  {\endgroup\par\addvspace{9pt}}
\setlength{\aboverulesep}{0.5pt}\setlength{\belowrulesep}{0.6pt}
\arrayrulecolor{filet}
\setlength{\heavyrulewidth}{0.9pt}\setlength{\lightrulewidth}{0.45pt}

% ---------------------------------------------------------------- encadrés
\newtcolorbox{callout}{enhanced,breakable,
  colback=fondclair,colframe=accent,boxrule=0pt,leftrule=2.6pt,
  sharp corners,left=11pt,right=11pt,top=8pt,bottom=8pt,
  before skip=10pt,after skip=11pt,parbox=false}
\newtcolorbox{codebox}{enhanced,breakable,
  colback=fondcode,colframe=filet,boxrule=0.5pt,arc=2pt,
  left=9pt,right=9pt,top=7pt,bottom=7pt,
  before skip=10pt,after skip=11pt}

% ---------------------------------------------------------------- sommaire
\renewcommand{\contentsname}{Sommaire}
\titlecontents{section}[0pt]
  {\addvspace{4pt}\sffamily\bfseries\color{encre}}
  {\contentslabel[\thecontentslabel]{1.6em}}{}
  {\titlerule*[0.7pc]{\color{filet}.}\contentspage}
\titlecontents{subsection}[1.9em]
  {\sffamily\small\color{gris}}
  {\contentslabel[\thecontentslabel]{2.3em}}{}
  {\titlerule*[0.7pc]{\color{filet}.}\contentspage}

% ---------------------------------------------------------------- divers
\newlength{\metalab}
\newlength{\tw}\setlength{\tw}{\textwidth}
\newcommand{\settw}[1]{\setlength{\tw}{\dimexpr\textwidth-#1pt\relax}}
% tolérance relâchée : sans elle, TeX préfère déborder de quelques points
% plutôt que de desserrer une ligne contenant un long jeton en chasse fixe.
\emergencystretch=1.6em
